\SourcePage{13}{16}
the meaning of ``with sufficient accuracy'' depends on the specific problem
under consideration. In a Wilson chamber, for instance, one tracks the
electron by its mist trail, whose thickness far exceeds atomic dimensions;
at such a level of accuracy in locating the trajectory, the electron is a
fully classical object.

Quantum mechanics thus occupies a very distinctive place among physical
theories: it contains classical mechanics as its limiting case and, at the
same time, requires that limiting case for its own foundation.

We can now formulate the problem addressed by quantum mechanics. A typical
problem is to predict the outcome of a repeated measurement from the known
outcomes of earlier measurements. We shall also see below that, compared with
classical mechanics, quantum mechanics generally restricts the set of values
that various physical quantities (energy, for example) may assume, that is,
the values that may be found by measuring the quantity in question. The
formalism of quantum mechanics must make it possible to determine these
allowed values.

In quantum mechanics the measurement process has an essential peculiarity: it
always acts upon the electron being measured, and for a given measurement
accuracy this action cannot in principle be made arbitrarily weak. The more
accurate the measurement, the stronger its action; only in measurements of
very low accuracy can the action on the measured object be weak. This property
of measurements is logically connected with the fact that the dynamical
characteristics of an electron arise only through the measurement itself. It
is clear that, if the action of the measurement process on the object could be
made arbitrarily weak, the quantity being measured would have a definite value
by itself, irrespective of the measurement.

Of the various types of measurement, the principal one is the measurement of
an electron's coordinates. Within the domain of applicability of
quantum mechanics, the electron's coordinates can always be
measured\footnote{We emphasize once again that by a ``measurement made'' we
mean the interaction of the electron with a classical ``apparatus''; this in
no way presupposes the presence of an outside observer.} to any accuracy.

Suppose that successive measurements of the coordinates%
\SourcePage{14}{17}
of an electron are made at specified time intervals $\Delta t$. As a rule,
these results will fail to follow any smooth curve. Quite the opposite:
the greater the measurement precision, the more abrupt and erratic
the succession of results becomes, consistent with the absence of a path
concept for the electron. A more or less smooth path is obtained only when the
electron's coordinates are measured with a low degree of accuracy, for
example, by the condensation of droplets of vapour in a Wilson chamber.

If instead, while keeping the measurement accuracy unchanged, the intervals
$\Delta t$ between measurements are reduced, adjacent measurements will, of
course, give nearby coordinate values. Yet although the results of a series of
successive measurements will lie in a small region of space, within that
region they will be distributed quite irregularly and will by no means fit any
smooth curve. In particular, as $\Delta t$ tends to zero, the results of
nearby measurements do not tend at all to lie on one straight line.

This last observation makes plain that a classical particle velocity does not
exist in quantum mechanics: the ratio of the coordinate change between two
time instants to the separation $\Delta t$ tends to no well-defined limit. As we shall see later, however, quantum mechanics
does admit a meaningful definition of a particle's velocity at a given
instant, one that reduces to the classical velocity upon passing to classical
mechanics.

In classical mechanics a particle possesses, at each instant, definite
coordinates and velocity; in quantum mechanics the situation is entirely
different. If a measurement has given an electron definite coordinates, it
then has no definite velocity at all. Conversely, an electron having a
definite velocity cannot have a definite position in space. Indeed, the
simultaneous existence of coordinates and velocity at any instant would imply
the presence of a definite velocity, which the electron does not have. Thus,
in quantum mechanics the electron's coordinates and velocity are quantities
that cannot be measured exactly at the same time, that is, cannot
simultaneously have definite values. One may say that the electron's
coordinates and velocity are quantities which do not exist simultaneously. A
quantitative relation governing the possibility of measuring coordinates and
velocity inexactly at the same instant will be derived below.
\SourcePage{15}{18}
In classical mechanics, the state of a physical system is completely described
by specifying all its coordinates and velocities at a given instant; from
these initial data the equations of motion fully determine the system's
behaviour at all future instants. In quantum mechanics
such a description is impossible in principle, since the coordinates and
their corresponding velocities do not exist simultaneously. Hence the state
of a quantum system is described by fewer quantities than in classical
mechanics, that is, its description is less detailed than a classical one.

From this follows a very important consequence regarding the character of
predictions in quantum mechanics. Whereas a classical description is
sufficient to predict with complete precision the future motion of a
mechanical system, the less detailed quantum-mechanical description manifestly
cannot be adequate for this purpose.
Consequently, even when an electron is in a state described as completely as
possible, its behaviour at subsequent instants is in principle not unique.
Quantum mechanics therefore cannot make strictly definite predictions about
the electron's future behaviour. For a specified initial state of an electron,
a subsequent measurement may give different results. The task of quantum
mechanics is only to determine the probability that one or another result will
be obtained in that measurement. In some cases, of course, the probability of
a particular measurement result may equal unity, that is, become a certainty,
so that the result of the measurement is unique.

Every measurement process in quantum mechanics falls into one of two
categories. The first, which includes the majority of measurements, comprises
those that yield no determinate result with certainty for any state of the
system. The second comprises measurements such that for every possible result
there exists a state in which that result is obtained with certainty. It is
precisely these latter measurements, which may be called
\emph{predictable}, that have the principal role in quantum mechanics. The
quantitative characteristics of a state determined by such measurements are
what quantum mechanics calls physical quantities. If in some state a
measurement gives a unique result with certainty, we shall say that the
corresponding physical quantity has a definite value in that state. Below we
shall everywhere%
\SourcePage{16}{19}
understand the expression ``physical quantity'' in precisely the sense stated
here.

We shall repeatedly find below that by no means every collection of physical
quantities in quantum mechanics can be measured simultaneously, that is, can
have definite values simultaneously (we have already discussed one example:
the velocity and coordinates of an electron).

Sets of physical quantities with the following property have an important
role in quantum mechanics: the quantities can be measured simultaneously and,
when they simultaneously have definite values, no other physical quantity
(unless it is a function of them) can have a definite value in that state. We
shall refer to such sets of physical quantities as \emph{complete sets}.

Every description of an electron's state arises from some measurement. We now
state what a complete description of a state means in quantum mechanics.
Completely described states arise from the simultaneous measurement of a
complete set of physical quantities. From the results of such a measurement
one can, in particular, determine the probabilities of the results of any
subsequent measurement, independently of everything that happened to the
electron before the first measurement.

Below, everywhere except in \secref{14}, by states of a quantum system we shall mean
states described in precisely this complete manner.
